\documentclass[11pt]{article}
\usepackage[margin=1in]{geometry}
\usepackage{amsmath,amssymb}

\title{A flat connection with nonabelian tangent holonomy\\[2pt]
  \large Disproof of Conjecture 00000003576}
\author{}
\date{}

\begin{document}
\maketitle

\noindent
An open subset of the real Banach space $\mathbb C^3\cong\mathbb R^6$
admits a flat tangent connection whose linear holonomy group is $S_3$.

\section*{A covering and a flat connection}

Let $U\subset\mathbb C^3$ be the open set of ordered triples of distinct
points. For $z=(z_0,z_1,z_2)\in U$ define $q(z)=(a_2,a_1,a_0)$ by
\[
  \prod_{i=0}^2 (T-z_i)=T^3+a_2T^2+a_1T+a_0.
\]
The fibers of $q$ consist exactly of the triples obtained by permuting
the three coordinates. For $h\in\mathbb C^3$, the derivative $Dq_z(h)$
is the coefficient vector of
\[
 -\sum_{i=0}^2 h_i\prod_{j\ne i}(T-z_j).
\]
If this polynomial vanishes identically, evaluation at $T=z_i$ gives
$h_i\prod_{j\ne i}(z_i-z_j)=0$ for each $i$. The roots are distinct, so
$h=0$. Thus $Dq_z$ is invertible and $q$ is a local diffeomorphism. In
particular, its image $M=q(U)$ is open in $\mathbb R^6$.

Choose a small neighborhood $V$ of any $z\in U$ on which $q$ is a
diffeomorphism onto an open set. It may be chosen so that its six
permuted copies are disjoint, since the action is free. Because the
fibers are permutation orbits, the inverse image of $q(V)$ is the
disjoint union of those copies. Hence $q\colon U\to M$ is a covering
with deck action by permutations of the coordinate blocks.

Equip $U$ with the standard flat connection in its global coordinates.
The deck transformations are restrictions of linear permutations and
preserve this connection. In the coordinates of local inverse sheets
of $q$, put the connection coefficients on $M$ equal to zero. Any two
such coordinates differ locally by a constant linear permutation, so
these local connections agree on overlaps. The resulting tangent
connection on $M$ is flat and $q$ pulls it back to the standard one
on $U$.

\section*{Two lifted loops and their transport}

Take $x=((-2,0),(0,0),(2,0))\in U$, writing each complex coordinate as
a plane point. To interchange the first two or the last two ordered
points, let $c=(-1,0)$ or $c=(1,0)$ be the midpoint of the selected pair.
Replace the pair by $c+r,c-r$, keeping the third point fixed, and follow
these two successive segments with $0\le s\le1$ on each:
\[
 r=(s-1,s),\qquad r=(s,1-s).
\]
The segments meet at $r=(0,1)$. The vector $r$ never vanishes, so the
two moving points remain distinct. Their first coordinates lie in
$[-2,0]$ for the first swap and $[0,2]$ for the second, whereas the
fixed point has first coordinate $2$ or $-2$, respectively. The paths
therefore remain in $U$. They start at $x$ and end at $g_{01}x$ and
$g_{12}x$, where $g_{01}$ and $g_{12}$ are the respective adjacent
coordinate transpositions. Their projections under $q$ are loops
based at $q(x)$.

Upstairs parallel transport keeps a coordinate vector $v$ fixed.
Consequently a loop whose lift ends at $gx$ transports $Dq_x(v)$ to
$Dq_{gx}(v)$. Differentiating $q\circ g=q$ yields
\[
 Dq_{gx}\circ g=Dq_x,
 \qquad (Dq_x)^{-1}\circ Dq_{gx}=g^{-1}.
\]
Under the tangent identification by $Dq_x$, the two displayed loops
act by $g_{01}$ and $g_{12}$, since transpositions are involutions. If
$e_i$ denotes the vector with $(1,0)$ in block $i$ and zeros elsewhere,
then
\[
  g_{01}g_{12}e_0=e_1,\qquad g_{12}g_{01}e_0=e_2.
\]
Thus the linear holonomy is nonabelian. In fact, every loop lifts to a
path ending at a permutation of $x$, and the two displayed transport
operators generate the faithful block action of $S_3$.

\section*{Lean verification}

The Lean 4 project checks both piecewise affine paths at every rational
parameter $s/D$ with $D>0$ and $0\le s\le D$, using integer coordinate
numerators. It verifies collision freedom, segment gluing, and the
permuted endpoints. The formulas above establish the same claims over
the real intervals. The theorem
\texttt{C3576.lifted\_transports\_do\_not\_commute} checks the algebraic
transport deduction from the derivative identity, including the explicit
noncommuting witness. Run \texttt{lake build} in the \texttt{lean}
directory to verify the project.

\end{document}
